% GENERATED FILE. Edit canonical lessons, then run npm run generate:books.
% canonical-source-hash: fnv1a64:717ac9a010353031
% canonical-lessons: FR-C268-tout-d-abord, FR-C268-en-premier-lieu, FR-C268-par-la-suite, FR-C268-en-dernier-lieu, FR-C268-pour-finir

\chapter{Tout d'abord}
\label{ch:fr-tout-d-abord}
\hlchaptermodality{268}

\begin{chapteropening}
\textbf{By the end of this chapter:} \emph{I can put the steps of a story in order: first of all, in the first place, afterwards, lastly and to finish.}

Complete the last lesson of chapter 268: I can put the steps of a story in order: first of all, in the first place, afterwards, lastly and to finish.
\end{chapteropening}

\section[tout d'abord]{tout d'abord — first of all, to begin with}

\label{lesson:FR-C268-tout-d-abord}



\begin{quote}
Before the new one: say the French for a team, then the French for a birthday.
\end{quote}

\subsection*{You'll want to know: tout d'abord}
\textbf{tout d'abord} — \textquotedblleft{}first of all, to begin with\textquotedblright{}.

\begin{quote}
Tout d'abord, nous avons pris un café.
\end{quote}

\begin{quote}
\emph{First of all, we had a coffee.}
\end{quote}

A little stronger than \emph{d'abord} alone; it opens a list or a story.

\begin{cousinweb}
\emph{tout d'abord} joins \emph{tout} \textquotedblleft{}all\textquotedblright{}, inherited from Latin \emph{tōtus} \textquotedblleft{}whole\textquotedblright{} (Lewis \& Short), and \emph{d'abord} \textquotedblleft{}at first\textquotedblright{}, which is built on \emph{abord} \textquotedblleft{}approach\textquotedblright{}, from \emph{aborder} \textquotedblleft{}to come alongside\textquotedblright{}, itself built on \emph{bord} \textquotedblleft{}edge, side of a ship\textquotedblright{}.
\end{cousinweb}

\begin{grammarlens}[title={What you've built}]
One more word for telling, explaining and describing at B1.
\end{grammarlens}

\subsection*{Your turn}
\begin{itemize}
\raggedright
  \item \emph{Say it:} \emph{tout d'abord}
  \item \emph{Say it:} \emph{tout d'abord}, once more
  \item \emph{Say it:} a sentence of your own with \emph{tout d'abord}
  \item \emph{Recall:} say the French for a team, then the French for a birthday, then say \emph{tout d'abord} again
\end{itemize}

\begin{tcolorbox}[breakable,colback=teal!4,colframe=teal!35!black,title={Before you move on}]
What does tout d'abord mean? (\textquotedblleft{}First of all, to begin with\textquotedblright{}.) Say it once more.
\end{tcolorbox}
\section[en premier lieu]{en premier lieu — in the first place, firstly}

\label{lesson:FR-C268-en-premier-lieu}



\begin{quote}
Before the new one: say the French for a birthday, then the French for first of all.
\end{quote}

\subsection*{You'll want to know: en premier lieu}
\textbf{en premier lieu} — \textquotedblleft{}in the first place, firstly\textquotedblright{}.

\begin{quote}
En premier lieu, il a parlé de sa famille.
\end{quote}

\begin{quote}
\emph{In the first place, he talked about his family.}
\end{quote}

Formal; it orders the points of a speech or a report, often followed by \emph{en second lieu}.

\begin{cousinweb}
\emph{lieu} is inherited from Latin \emph{locus} \textquotedblleft{}place\textquotedblright{}, and \emph{premier} comes from Latin \emph{prīmārius} \textquotedblleft{}of the first rank\textquotedblright{}, built on \emph{prīmus} \textquotedblleft{}first\textquotedblright{} (EtymDB).
\end{cousinweb}

\begin{grammarlens}[title={What you've built}]
One more word for telling, explaining and describing at B1.
\end{grammarlens}

\subsection*{Your turn}
\begin{itemize}
\raggedright
  \item \emph{Say it:} \emph{en premier lieu}
  \item \emph{Say it:} \emph{en premier lieu}, once more
  \item \emph{Say it:} a sentence of your own with \emph{en premier lieu}
  \item \emph{Recall:} say the French for a birthday, then the French for first of all, then say \emph{en premier lieu} again
\end{itemize}

\begin{tcolorbox}[breakable,colback=teal!4,colframe=teal!35!black,title={Before you move on}]
What does en premier lieu mean? (\textquotedblleft{}In the first place, firstly\textquotedblright{}.) Say it once more.
\end{tcolorbox}
\section[par la suite]{par la suite — afterwards, later on}

\label{lesson:FR-C268-par-la-suite}



\begin{quote}
Before the new one: say the French for first of all, then the French for in the first place.
\end{quote}

\subsection*{You'll want to know: par la suite}
\textbf{par la suite} — \textquotedblleft{}afterwards, later on\textquotedblright{}.

\begin{quote}
Par la suite, ils sont devenus amis.
\end{quote}

\begin{quote}
\emph{Afterwards, they became friends.}
\end{quote}

It looks back from a later point in the story: what came next.

\begin{cousinweb}
\emph{suite} is inherited from Vulgar Latin \emph{sequita} \textquotedblleft{}a following\textquotedblright{} (EtymDB), built on Latin \emph{sequī} \textquotedblleft{}to follow\textquotedblright{}.
\end{cousinweb}

\begin{grammarlens}[title={What you've built}]
One more word for telling, explaining and describing at B1.
\end{grammarlens}

\subsection*{Your turn}
\begin{itemize}
\raggedright
  \item \emph{Say it:} \emph{par la suite}
  \item \emph{Say it:} \emph{par la suite}, once more
  \item \emph{Say it:} a sentence of your own with \emph{par la suite}
  \item \emph{Recall:} say the French for first of all, then the French for in the first place, then say \emph{par la suite} again
\end{itemize}

\begin{tcolorbox}[breakable,colback=teal!4,colframe=teal!35!black,title={Before you move on}]
What does par la suite mean? (\textquotedblleft{}Afterwards, later on\textquotedblright{}.) Say it once more.
\end{tcolorbox}
\section[en dernier lieu]{en dernier lieu — lastly, last of all}

\label{lesson:FR-C268-en-dernier-lieu}



\begin{quote}
Before the new one: say the French for in the first place, then the French for afterwards.
\end{quote}

\subsection*{You'll want to know: en dernier lieu}
\textbf{en dernier lieu} — \textquotedblleft{}lastly, last of all\textquotedblright{}.

\begin{quote}
En dernier lieu, elle a fermé la porte à clé.
\end{quote}

\begin{quote}
\emph{Last of all, she locked the door.}
\end{quote}

The partner of \emph{en premier lieu}; careful, written French.

\begin{cousinweb}
\emph{dernier} comes from Old French \emph{derrenier}, usually traced to Latin \emph{dē retrō} \textquotedblleft{}from behind\textquotedblright{} (Lewis \& Short); \emph{lieu} is inherited from Latin \emph{locus} \textquotedblleft{}place\textquotedblright{}.
\end{cousinweb}

\begin{grammarlens}[title={What you've built}]
One more word for telling, explaining and describing at B1.
\end{grammarlens}

\subsection*{Your turn}
\begin{itemize}
\raggedright
  \item \emph{Say it:} \emph{en dernier lieu}
  \item \emph{Say it:} \emph{en dernier lieu}, once more
  \item \emph{Say it:} a sentence of your own with \emph{en dernier lieu}
  \item \emph{Recall:} say the French for in the first place, then the French for afterwards, then say \emph{en dernier lieu} again
\end{itemize}

\begin{tcolorbox}[breakable,colback=teal!4,colframe=teal!35!black,title={Before you move on}]
What does en dernier lieu mean? (\textquotedblleft{}Lastly, last of all\textquotedblright{}.) Say it once more.
\end{tcolorbox}
\section[pour finir]{pour finir — to finish, to end with}

\label{lesson:FR-C268-pour-finir}



\begin{quote}
Before the new one: say the French for afterwards, then the French for lastly.
\end{quote}

\subsection*{You'll want to know: pour finir}
\textbf{pour finir} — \textquotedblleft{}to finish, to end with\textquotedblright{}.

\begin{quote}
Pour finir, nous avons chanté tous ensemble.
\end{quote}

\begin{quote}
\emph{To finish, we all sang together.}
\end{quote}

Common at the end of a story or a meal: \emph{pour finir, un petit café}.

\begin{cousinweb}
\emph{finir} is inherited from Latin \emph{fīnīre} \textquotedblleft{}to bound, to bring to an end\textquotedblright{} (EtymDB), built on \emph{fīnis} \textquotedblleft{}boundary, end\textquotedblright{} (Lewis \& Short).
\end{cousinweb}

\begin{grammarlens}[title={What you've built}]
One more word for telling, explaining and describing at B1.
\end{grammarlens}

\subsection*{Your turn}
\begin{itemize}
\raggedright
  \item \emph{Say it:} \emph{pour finir}
  \item \emph{Say it:} \emph{pour finir}, once more
  \item \emph{Say it:} a sentence of your own with \emph{pour finir}
  \item \emph{Recall:} say the French for afterwards, then the French for lastly, then say \emph{pour finir} again
\end{itemize}

\begin{tcolorbox}[breakable,colback=teal!4,colframe=teal!35!black,title={Before you move on}]
What does pour finir mean? (\textquotedblleft{}To finish, to end with\textquotedblright{}.) Say it once more.
\end{tcolorbox}
